\section*{Chapitre \ref{ch:ml:experiment-tracking-and-reproducibility} --- Suivi des expériences et reproductibilité}

\begin{solution}{exo:ml:experiment-tracking-and-reproducibility:1}
Paramètres (une valeur par défaut change d'une version de bibliothèque à l'autre), graine (l'échantillon de lignes et de colonnes que voit chaque arbre), instantané de données
(une révision du fournisseur), version du code (la définition d'une variable), environnement (versions des bibliothèques, nombres de threads, matériel
modifiant l'ordre des sommes en virgule flottante).
\end{solution}

\begin{solution}{exo:ml:experiment-tracking-and-reproducibility:2}
Une empreinte (hash) du contenu identifie l'exécution par ce qu'elle est~: deux fiches identiques reçoivent un même identifiant, une fiche ne peut être
modifiée sans changer son identifiant, et rien ne dépend du moment ou du lieu d'écriture, si bien que l'outil de suivi lui-même est reproductible.
\end{solution}

\begin{solution}{exo:ml:experiment-tracking-and-reproducibility:3}
Le ratio de Sharpe déflaté est la probabilité que le vrai ratio de Sharpe du champion dépasse le meilleur résultat que 200
essais sans compétence seraient censés afficher~: 0.15, une preuve faible. Le ratio de Sharpe probabiliste le compare à zéro,
comme s'il était le seul essai~: 0.96. L'écart est le prix de la recherche.
\end{solution}

\begin{solution}{exo:ml:experiment-tracking-and-reproducibility:4}
La formule suppose 200 essais indépendants~; les exécutions partagent les données et ne diffèrent que par quelques paramètres, si bien que leurs rendements en validation sont
fortement corrélés, et le nombre effectif d'essais indépendants est bien plus petit. Le maximum attendu de moins d'essais est
plus faible, et le ratio de Sharpe déflaté plus élevé~; le nombre effectif peut être estimé à partir des corrélations entre exécutions.
\end{solution}

\begin{solution}{exo:ml:experiment-tracking-and-reproducibility:5}
Chaque entrée contient l'empreinte de l'entrée précédente, et sa propre empreinte couvre son contenu~: modifier ou supprimer une ligne change son empreinte
et rompt le lien de toutes les lignes suivantes. Cela n'empêche pas de réécrire toute la chaîne à partir de la ligne modifiée~; pour
cela, la dernière empreinte doit être ancrée dans un lieu inaccessible à celui qui modifie (un système séparé, un condensat quotidien signé).
\end{solution}

\begin{solution}{exo:ml:experiment-tracking-and-reproducibility:6}
Les poids permettent d'exécuter le modèle, non de le reconstruire~: sans l'instantané de données, le code, les paramètres et la graine, on ne peut
ni le réentraîner, ni auditer la façon dont il a été choisi, ni produire une version corrigée qui ne diffère que par la correction. Dans le chapitre, chaque champ manquant a modifié
le ratio de Sharpe de test du modèle reconstruit d'un écart de 0.44 à 1.28.
\end{solution}

\begin{solution}{exo:ml:experiment-tracking-and-reproducibility:7}
Dix graines donnent des ratios de Sharpe de test allant de $-0.20$ à 2.88, de moyenne 1.03 et d'écart type 0.94. La configuration du champion est
banale une fois sa graine chanceuse moyennée~: son ratio de Sharpe de test de 0.91 est proche de la moyenne, et la dispersion est aussi grande que
la moyenne.
\end{solution}

\begin{solution}{exo:ml:experiment-tracking-and-reproducibility:8}
Pour chaque version hebdomadaire~: la \omterm{def:ml:experiment-tracking-and-reproducibility:tracker}{fiche d'exécution} (instantané de données, version du code, paramètres, graine, métriques), qui a déclenché le réentraînement (le
planificateur, sous l'autorité de qui), qui a approuvé la mise en production et sur quelle base, la nature du changement (nouvelles données seulement, ou nouveau
code ou paramètres), les dates d'enregistrement, de mise en production et de retrait~; la \omterm{def:ml:experiment-tracking-and-reproducibility:registry}{piste d'audit} ancrée et conservée pendant la durée que
les règles exigent.
\end{solution}

\begin{solution}{pb:ml:experiment-tracking-and-reproducibility:1}
\textbf{Partie I.}
\begin{enumerate}
\item Paramètres, graine, instantané de données, version du code, environnement, métriques, empreinte de l'artefact et horodatage~; l'identifiant est
l'empreinte de la fiche.
\item De $-3.16$ à 2.22, médiane $-0.83$~; validation du champion 2.22 et test 0.91.
\item 0.96 et 0.15.
\item 3.58, au-dessus des 2.22 du champion~: le champion n'est pas distinguable de la chance de 200 tentatives indépendantes (l'indice de référence
est conservateur pour des exécutions corrélées).
\end{enumerate}
\textbf{Partie II.}
\begin{enumerate}[resume]
\item Le fournisseur a révisé 1~\% des rendements (instantané 2026-03), et l'écrêtage des variables est passé de 3 à 2.5 écarts types (v2).
\item D'après l'historique des étapes du registre, à la date du 10 mars à 08:00~: version 1, mise en production le 2 mars.
\item Enregistrement et mise en production avec dates et motifs, chaque entrée liée par empreinte à la précédente~; modifier une ligne rompt la chaîne.
\item Fiches de chaque modification substantielle d'un logiciel de trading algorithmique (quand, par qui, approuvée par qui, en quoi)~; registres d'ordres des
traders HFT, conservés cinq ans.
\end{enumerate}
\textbf{Partie III.}
\begin{enumerate}[resume]
\item Le vidage du modèle reconstruit a l'empreinte du contenu de l'artefact stocké~: le même modèle, bit à bit.
\item Chaque champ manquant rompt la reproduction~: paramètres par défaut $-1.28$, graine $+0.85$, dernières données $-1.25$, code actuel
$-0.44$ sur le ratio de Sharpe de test.
\item Ce seul bruit de graine modifie le résultat autant que le choix de la configuration.
\item Entre versions de bibliothèques, nombres de threads ou matériels, là où l'ordre des opérations en virgule flottante change~; là, la norme est
d'obtenir les mêmes métriques à l'intérieur de la dispersion due à la graine.
\end{enumerate}
\textbf{Partie IV.}
\begin{enumerate}[resume]
\item \emph{Reproduire le modèle de mars.} Deux cents exécutions se cachent derrière le champion, dont le ratio de Sharpe en validation de 2.22 a un
ratio de Sharpe déflaté de 0.15~; la reproduction depuis le registre est exacte, et elle échoue dès qu'il manque un seul élément parmi les paramètres, la graine, l'instantané de données
ou la version du code.
\item Pas sur ces seules preuves~: il faut réclamer la dispersion due aux graines, le nombre effectif d'essais et une validation plus longue.
\item Interroger le registre à la date voulue, remonter le lignage jusqu'à l'exécution, la reconstruire et lister les 200 essais depuis l'outil de suivi~:
des minutes, non des semaines.
\item Conservation et immutabilité des fiches et des artefacts pendant la période, et ancrage de la \omterm{def:ml:experiment-tracking-and-reproducibility:registry}{piste d'audit}.
\item Comme borne inférieure~; le nombre réel est plus élevé, et le ratio de Sharpe déflaté est une borne supérieure.
\item Chaque réentraînement comme une version avec son instantané, une règle de mise en production automatique enregistrée comme approbateur, et des alertes lorsque la
règle est contournée.
\item Une personne autre que l'auteur du modèle, selon la \omterm{def:ml:reinforcement-learning-foundations:mdp}{politique} de risque de modèle de la société (livre~6, chapitre~26).
\item La fiche complète de l'exécution et le même logiciel.
\end{enumerate}
\end{solution}

\begin{solution}{iq:ml:experiment-tracking-and-reproducibility:1}
Paramètres, graines, versions des données et du code, environnement, métriques et artefacts, pour chaque exécution. Les exécutions en échec comptent comme essais~: sans
elles, la correction des tests multiples est impossible et le résultat retenu paraît meilleur qu'il ne l'est.
\iqlookfor{les champs, et les exécutions échouées comptées comme essais.}
\end{solution}

\begin{solution}{iq:ml:experiment-tracking-and-reproducibility:2}
Fixer et consigner les graines, utiliser des algorithmes déterministes et des réductions mono-thread ou à ordre fixe, figer les versions des bibliothèques, versionner
les données et le code, et vérifier en reconstruisant et en comparant les empreintes de contenu.
\iqlookfor{réglages de déterminisme, versions figées et comparaison d'empreintes.}
\end{solution}

\begin{solution}{iq:ml:experiment-tracking-and-reproducibility:3}
Le meilleur de nombreux essais est biaisé à la hausse~; le ratio de Sharpe déflaté le compare au maximum attendu d'autant d'essais
nuls, si bien que leur nombre (et leur corrélation) détermine à quel point un ratio de Sharpe est impressionnant.
\iqlookfor{biais de sélection et ratio de Sharpe déflaté.}
\end{solution}

\begin{solution}{iq:ml:experiment-tracking-and-reproducibility:4}
Nom, version, étape (candidat, production, retiré), lignage (exécution, données, code), métriques, propriétaire et approbateur, historique des étapes avec
les motifs, les liens vers la \omterm{def:ml:interpretability-and-model-governance:card}{fiche de modèle} et le rapport de validation, et une \omterm{def:ml:experiment-tracking-and-reproducibility:registry}{piste d'audit} en ajout seul des modifications.
\iqlookfor{étapes, lignage, approbations et historique.}
\end{solution}

\begin{solution}{iq:ml:experiment-tracking-and-reproducibility:5}
Parce que le même nom finit par désigner des données différentes~; les corrections créent de nouveaux instantanés immuables (indexés par empreinte de contenu ou datés) et les exécutions
enregistrent ceux qu'elles ont utilisés, si bien que les anciens résultats restent reproductibles et que les nouveaux peuvent être comparés.
\iqlookfor{immutabilité et identifiants d'instantanés.}
\end{solution}

\begin{solution}{iq:ml:experiment-tracking-and-reproducibility:6}
Partir du lignage du registre~: la \omterm{def:ml:experiment-tracking-and-reproducibility:tracker}{fiche d'exécution} est-elle complète, l'instantané de données est-il encore disponible et inchangé (empreinte), la
version du code extraite, l'environnement identique~; reconstruire et comparer les empreintes~; comparer les sorties intermédiaires étape par étape (manifestes de pipeline
du livre~7) pour trouver l'endroit où elles divergent.
\iqlookfor{un parcours systématique, du lignage jusqu'à la première étape divergente.}
\end{solution}
